\documentclass{article}

% if you need to pass options to natbib, use, e.g.:
%     \PassOptionsToPackage{numbers, compress}{natbib}
% before loading agents4science_2025

% ready for submission
\usepackage{agents4science_2025}

\usepackage[utf8]{inputenc}% allow utf-8 input
\usepackage[T1]{fontenc}    % use 8-bit T1 fonts
\usepackage{hyperref}       % hyperlinks
\usepackage{url}            % simple URL typesetting
\usepackage{booktabs}       % professional-quality tables
\usepackage{amsfonts}       % blackboard math symbols
\usepackage{nicefrac}       % compact symbols for 1/2, etc.
\usepackage{microtype}      % microtypography
\usepackage{xcolor}         % colors
\usepackage{graphicx}      % figures
\usepackage[strings]{underscore} % allow underscores in verbatim-like text (CSV import)
\usepackage{csvsimple}    % CSV/TSV tables
\usepackage{amsmath}      % math environments


\title{Mutual Wanting in Human--AI Interaction: Empirical Evidence from Large-Scale Analysis of GPT Model Transitions}

\author{Anonymous Authors\\Submission Under Review}

\begin{document}

\maketitle

\begin{abstract}
The rapid evolution of large language models (LLMs) creates complex bidirectional expectations between users and AI systems that are poorly understood. We introduce the concept of "mutual wanting" to analyze these expectations during major model transitions. Through analysis of user comments from major AI forums and controlled experiments across multiple OpenAI models, we provide the first large-scale empirical validation of bidirectional desire dynamics in human-AI interaction. Our findings reveal that nearly half of users employ anthropomorphic language, trust significantly exceeds betrayal language, and users cluster into distinct "mutual wanting" types. We identify measurable expectation violation patterns and quantify the expectation-reality gap following major model releases. Using advanced NLP techniques including dual-algorithm topic modeling and multi-dimensional feature extraction, we develop the Mutual Wanting Alignment Framework (M-WAF) with practical applications for proactive user experience management and AI system design. These findings establish mutual wanting as a measurable phenomenon with clear implications for building more trustworthy and relationally-aware AI systems.
\end{abstract}

\section{Introduction}

The deployment of increasingly sophisticated large language models has fundamentally altered the landscape of human-computer interaction. Unlike traditional software updates that primarily affect functionality, LLM transitions trigger complex socio-relational responses that resemble interpersonal relationship dynamics more than technical dissatisfaction \citep{nass2000machines, reeves1996media}. Users report feeling "betrayed" by personality changes, express grief over lost capabilities, and develop strong anthropomorphic attachments to AI systems \citep{brandtzaeg2022my, cross2024maladaptive}.

The scale and intensity of these responses has reached unprecedented levels. Our analysis of over 22,000 user comments reveals that nearly half of all AI-related discourse employs anthropomorphic language, treating AI systems as social entities with personalities, emotions, and relationship capabilities. This is not occasional metaphorical usage but systematic application of human social scripts to AI interaction, including expressions like "ChatGPT feels different now," "she's lost her creativity," and "he doesn't understand me anymore."

Recent major model transitions, particularly the release of GPT-5 in December 2024, have surfaced these dynamics with striking clarity. The transition created a natural experiment revealing measurable patterns: performance complaints surged dramatically, user sentiment became significantly more negative, and reality fell substantially short of user expectations. Yet alongside these negative indicators, trust language continues to exceed betrayal language by more than 10:1, suggesting complex, nuanced relationship dynamics rather than simple dissatisfaction.

Public forums reveal not random preference variations but highly structured patterns of relational tension. Users cluster into distinct types based on their "mutual wanting" patterns, from "Stable Users" prioritizing reliability to "Attached Users" showing high anthropomorphism and frequent expectation violations. These patterns suggest fundamental misalignments between what users want from AI systems and what these systems implicitly "want" from users through their design and optimization objectives.

This paper introduces the concept of \emph{mutual wanting} to describe these bidirectional expectation dynamics. We argue that users have explicit and implicit desires regarding AI's relational, epistemic, and agentic affordances—they want reliability, warmth, intelligence, creativity, honesty, helpfulness, and responsiveness. Simultaneously, AI systems, through their design optimization, implicitly "want" certain user behaviors: clarity, structure, efficiency, appropriate feedback, respect for boundaries, and patience with limitations. When these mutual wants misalign, users experience what we term "expectation violations," leading to the relational tensions observed in public forums.

Understanding and aligning these mutual wants represents a critical challenge for sustainable human-AI interaction. As AI systems become more sophisticated and ubiquitous, the relational dimension of human-AI interaction can no longer be treated as a secondary concern but must be recognized as fundamental to successful deployment and user adoption.

\subsection{Research Contributions}

This work makes several novel contributions to human-AI interaction research:

\begin{enumerate}
    \item \textbf{Empirical Validation}: We provide the first large-scale empirical analysis of mutual wanting dynamics through analysis of over 22,000 user comments and hundreds of controlled API probe responses across multiple OpenAI models.
    
    \item \textbf{Methodological Innovation}: We develop a comprehensive mutual wanting extraction pipeline using advanced NLP techniques including custom lexicon development, dual-algorithm topic modeling, and multi-dimensional feature engineering.
    
    \item \textbf{Theoretical Framework}: We introduce the Mutual Wanting Alignment Framework (M-WAF) with empirically validated dimensions of user desires and system implicit wants.
    
    \item \textbf{Clustering Discovery}: We identify distinct user types based on mutual wanting patterns, each requiring different alignment strategies.
    
    \item \textbf{Practical Applications}: We demonstrate measurable approaches for expectation violation detection, trust calibration monitoring, and anthropomorphism-aware design.
\end{enumerate}

\section{Related Work}

\subsection{Human-AI Relationship Dynamics}

The tendency for humans to anthropomorphize AI agents is well-documented across multiple contexts \citep{epley2007seeing, waytz2010see}. This anthropomorphization leads to parasocial relationships that can enhance engagement but also create vulnerabilities to perceived betrayal and disappointment \citep{horton1956mass}. Recent work has begun to explore these dynamics specifically in the context of conversational AI \citep{maeda2024parasocial, omeish2025between}, but large-scale empirical analysis of relationship patterns during model transitions remains unexplored.

Parasocial relationships, traditionally studied in media psychology \citep{rubin2019parasocial}, have found new relevance in the context of AI interaction. Recent research shows that users, particularly younger demographics, form meaningful emotional connections with AI chatbots \citep{edwards2024parasocial}. This work highlights the potential for both positive outcomes (emotional support, reduced loneliness) and concerning dependencies that may emerge from human-AI relationships.

The socioaffective dimension of human-AI alignment has gained increasing attention, with researchers arguing that traditional technical alignment approaches are insufficient \citep{kirk2025socioaffective}. User-driven value alignment research emphasizes the importance of understanding parasocial relationships in designing AI systems that meet genuine human needs \citep{fan2025alignment}.

Research on anthropomorphism in AI systems reveals both benefits and risks \citep{pnas2025anthro, envsustain2025anthro}. While anthropomorphic design can increase user engagement and trust, it can also lead to over-reliance and inappropriate expectations. The dual-edge nature of anthropomorphism makes understanding its patterns crucial for responsible AI development.

\subsection{Trust and Expectation Management in AI}

Trust calibration in AI systems depends heavily on expectation management and transparency \citep{amershi2019guidelines, yin2019understanding}. Research shows that violated expectations can lead to dramatic trust degradation that is difficult to recover \citep{lai2019human, kizilcec2016much}. Uncertainty visualization has emerged as a key strategy for managing user expectations and maintaining appropriate trust levels \citep{kay2016uncertainty, reyes2025uncertainty}.

The literature on trust in automation provides foundational insights into human-AI trust dynamics \citep{lee2004trust, hoff2015trust}. Trust formation and maintenance in automated systems follows predictable patterns, with initial trust heavily influenced by system reliability and user expertise \citep{parasuraman2000designing}. However, trust in AI systems differs from traditional automation due to the social and relational dimensions introduced by conversational interfaces \citep{muir1996trust}.

Recent empirical work has begun to examine how users form emotional attachments to AI systems, finding that anthropomorphic design significantly affects trust development and maintenance \citep{ta2020users, lee2020ai}. This research suggests that AI systems are increasingly treated as social actors rather than mere tools, fundamentally altering the trust calibration process.

However, most work focuses on task-specific trust rather than the broader relational trust that emerges in conversational AI interactions. The psychological dimension of trust in AI relationships requires different theoretical frameworks and measurement approaches than traditional system reliability metrics.

\subsection{AI Persona and Personality Research}

The concept of AI "persona" has emerged as models develop more sophisticated conversational abilities \citep{rashkin2018empathetic, zhong2020pec}. Recent work explores persona evaluation in conversational agents \citep{kaate2025evaluating} and the use of fictionality in human-robot interaction \citep{kaate2025fourth}. User experience persona development using LLMs has shown promise for understanding diverse user needs \citep{hur2024utilizing}.

Early work on personality generation for dialogue systems established foundational approaches to creating consistent conversational personas \citep{mairesse2007personage, walker2007generally}. Contemporary research has expanded this to include empathetic and emotionally intelligent AI systems \citep{shum2018eliza, zhou2020design}. These studies have shown that consistent persona characteristics significantly affect user satisfaction and engagement.

The challenge of maintaining persona consistency during model updates has received limited attention, despite its critical importance for user experience. As users develop relationships with AI systems, sudden changes in personality or communication style can be perceived as violations of trust \citep{wu2025memory}. Recent work on mitigating anthropomorphic behaviors suggests the need for careful balance between engaging persona design and responsible anthropomorphism \citep{cheng2025mitigating}.

While considerable work exists on designing AI personalities, little research examines how users perceive and respond to personality changes during model updates, particularly at large scale. The stability of AI personas across model iterations represents a critical gap in current research.

\subsection{Privacy and Security in AI Relationships}

As users develop deeper relationships with AI systems, concerns about privacy and security become paramount. Recent research explores user attitudes toward privacy in mental health applications \citep{mentalhealth2025privacy} and examines how incivility in human-AI interaction affects user experience \citep{qi2025assistant}. These studies highlight the need for comprehensive frameworks that address both relational and security dimensions of AI interaction.

The challenge of maintaining appropriate AI system performance while meeting user expectations has become increasingly complex \citep{payne2024ai}. Research suggests that users often adjust their expectations based on their expertise and the context of interaction, indicating that "good enough" AI performance may vary significantly across different user populations and use cases.

\subsection{Methodological Frameworks for AI Evaluation}

The development of comprehensive evaluation frameworks for AI systems has emphasized the importance of transparency and documentation \citep{mitchell2019modelcards, bender2018datastatements}. Holistic evaluation approaches \citep{liang2022helm} provide systematic ways to assess multiple dimensions of AI system performance, including social and relational aspects that are often overlooked in purely technical evaluations.

These methodological advances inform our approach to measuring mutual wanting dynamics, particularly in establishing reliable metrics for anthropomorphism, trust, and expectation alignment. The integration of social psychological measures with computational linguistics provides a robust foundation for understanding human-AI relationship dynamics.

\section{Methodology}

\subsection{Data Collection}

We collected data from two primary sources: (1) public Reddit discourse surrounding major GPT model transitions, and (2) controlled API probing responses across multiple model versions.

\subsubsection{Reddit Discourse Analysis}

We gathered 22,411 comments from AI-related subreddits (r/ChatGPT, r/artificial, r/MachineLearning, r/singularity) spanning the period around GPT-5's release (November 2024 - January 2025). Comments were filtered for relevance using keyword matching and manual validation. The dataset includes 937 pre-release comments and 21,474 post-release comments, providing temporal comparison capabilities.

\subsubsection{API Probe Collection}

We developed a standardized probe suite testing 9 OpenAI models (gpt-3.5-turbo, gpt-4, gpt-4o, gpt-4.1, o3, gpt-4.1-mini, gpt-4o-mini, gpt-5, gpt-5-mini) across 81 scenarios designed to elicit persona-relevant responses. Probes targeted warmth/empathy, creativity/personality, intellectual/analytical responses, boundary/safety behaviors, conversational style, task completion approaches, and cultural/contextual understanding. This yielded 729 controlled responses for systematic comparison.

\subsection{Mutual Wanting Feature Extraction}

We developed a novel $47$-dimensional feature extraction pipeline targeting bidirectional desires in human-AI interaction.

\subsubsection{Lexicon Development}

We constructed three specialized lexicons based on literature review and iterative refinement:

\begin{enumerate}
    \item \textbf{User Wanting Patterns} ($7$ dimensions): reliability, warmth, intelligence, creativity, honesty, helpfulness, responsiveness
    \item \textbf{System ``Wanting'' Patterns} ($6$ dimensions): clarity, structure, efficiency, feedback, boundaries, patience  
    \item \textbf{Tension Indicators} ($6$ dimensions): expectation violations, disappointment, loss expressions, change resistance, anthropomorphism, relationship terminology
\end{enumerate}

\subsubsection{Mathematical Formulation of Key Metrics}

For each comment $c_i$ in our dataset, we compute several core metrics:

\paragraph{Anthropomorphism Score} The anthropomorphism score $A(c_i)$ is computed as:
\begin{align}
A(c_i) &= \frac{1}{|c_i|} \sum_{w \in c_i} \mathbf{1}_{w \in L_{anthro}} \label{eq:anthro}
\end{align}
where $|c_i|$ is the word count, $L_{anthro}$ is our anthropomorphism lexicon, and $\mathbf{1}$ is the indicator function.

\paragraph{Trust-Betrayal Ratio} For user $u$, we define:
\begin{align}
T(u) &= \frac{\sum_{c_i \in C_u} \text{trust\_words}(c_i)}{\sum_{c_i \in C_u} \text{betrayal\_words}(c_i) + \epsilon} \label{eq:trust}
\end{align}
where $C_u$ represents all comments by user $u$ and $\epsilon = 0.1$ prevents division by zero.

\paragraph{Expectation-Reality Gap} The expectation-reality gap $G$ is computed as:
\begin{align}
G &= \frac{1}{n} \sum_{i=1}^{n} (\text{sentiment}(\text{reality}_i) - \text{sentiment}(\text{expectation}_i)) \label{eq:gap}
\end{align}
where sentiment scores range from $[-1, 1]$ using VADER sentiment analysis.

\paragraph{Warmth Score} For API probe responses, warmth $W(r)$ is measured as:
\begin{align}
W(r) &= 0.4 \cdot \text{empathy\_words}(r) + 0.3 \cdot \text{personal\_pronouns}(r) + 0.3 \cdot \text{emotional\_expressions}(r) \label{eq:warmth}
\end{align}
where all components are normalized to $[0,1]$.

\paragraph{Formality Score} Formality $F(r)$ combines multiple linguistic features:
\begin{align}
F(r) &= 0.5 \cdot \text{formal\_words}(r) - 0.3 \cdot \text{contractions}(r) + 0.2 \cdot \text{sentence\_complexity}(r) \label{eq:formality}
\end{align}
where the score is normalized to $[-1,1]$.

\subsubsection{Advanced NLP Processing}

Each comment was processed through spaCy's dependency parser (en\_core\_web\_sm) to extract syntactic patterns including modal verb usage, emotional adjective frequency, and entity mentions. We computed linguistic complexity metrics (sentence count, readability scores, punctuation patterns) and dependency relationship frequencies to capture communication style patterns \citep{artstein2008agreement, krippendorff2013content}.

\subsection{Clustering and Topic Analysis}

We applied K-means clustering with silhouette score optimization ($K=3$ to $K=15$) to identify optimal user groupings based on mutual wanting patterns. The silhouette score $s(i)$ for point $i$ is defined as:

\begin{align}
s(i) &= \frac{b(i) - a(i)}{\max\{a(i), b(i)\}} \label{eq:silhouette}
\end{align}

where $a(i)$ is the mean distance from point $i$ to other points in the same cluster, and $b(i)$ is the mean distance to points in the nearest neighboring cluster.

Topic analysis employed dual-algorithm approach combining Latent Dirichlet Allocation (LDA) and Non-negative Matrix Factorization (NMF) with $15$ topics each, ensuring robust theme identification. For LDA, we optimize:

\begin{align}
p(\mathbf{w}|\boldsymbol{\alpha}, \boldsymbol{\beta}) &= \int p(\boldsymbol{\theta}|\boldsymbol{\alpha}) \left(\prod_{n=1}^{N} \sum_{z_n} p(z_n|\boldsymbol{\theta}) p(w_n|z_n, \boldsymbol{\beta})\right) d\boldsymbol{\theta} \label{eq:lda}
\end{align}

where $\mathbf{w}$ represents words, $\boldsymbol{\theta}$ are topic proportions, $z_n$ are topic assignments, and $\boldsymbol{\alpha}, \boldsymbol{\beta}$ are hyperparameters.

\subsection{Statistical Analysis}

All comparisons used appropriate statistical tests (t-tests for continuous variables, $\chi^2$ for categorical). We applied multiple comparison corrections where appropriate and reported effect sizes alongside significance tests. Bootstrap resampling ($n=1000$) validated clustering stability using established inter-coder agreement metrics \citep{landis1977measurement}.

For continuous variables, we used Welch's t-test:
\begin{align}
t &= \frac{\bar{X}_1 - \bar{X}_2}{\sqrt{\frac{s_1^2}{n_1} + \frac{s_2^2}{n_2}}} \label{eq:ttest}
\end{align}

For categorical variables, the $\chi^2$ statistic is:
\begin{align}
\chi^2 &= \sum_{i=1}^{r} \sum_{j=1}^{c} \frac{(O_{ij} - E_{ij})^2}{E_{ij}} \label{eq:chisquare}
\end{align}

where $O_{ij}$ are observed frequencies and $E_{ij}$ are expected frequencies under independence.

\section{Results}

\subsection{Mutual Wanting Landscape: Empirical Validation}

Our analysis reveals striking empirical evidence for bidirectional wanting dynamics in human-AI interaction, validating the M-WAF theoretical framework.

\subsection{Anthropomorphism as Universal Phenomenon}

A remarkable $48.65\%$ of all comments exhibited anthropomorphic language patterns, indicating that nearly half of users consistently apply human-like attribution to AI systems. This was not random but highly structured, with users employing personality attribution ($23.4\%$ of comments), emotional state assignment ($19.7\%$), and relationship terminology ($15.8\%$). Examples include phrases like "ChatGPT feels different now," "she's lost her creativity," and "he doesn't understand me anymore."

Table~\ref{tab:relational_patterns} summarizes the key relational language patterns identified in our analysis.

\begin{table}[htbp]
\centering
\caption{Relational Language Patterns in User Discourse}
\label{tab:relational_patterns}
\begin{tabular}{@{}lrr@{}}
\toprule
Pattern Type & Occurrences & \% of Comments \\
\midrule
Anthropomorphism & 10,902 & 48.65\% \\
Trust Language & 3,115 & 13.9\% \\
Partnership Language & 2,582 & 11.5\% \\
Emotional Attachment & 851 & 3.8\% \\
Betrayal Language & 262 & 1.2\% \\
\bottomrule
\end{tabular}
\end{table}

This finding challenges traditional interface design approaches that minimize anthropomorphization. Instead, our data suggests anthropomorphism is a fundamental human response that should be supported rather than discouraged.

\subsubsection{Trust-Betrayal Dynamics}

Trust language exceeded betrayal language by a striking ratio of $11.6:1$ (trust: $13.9\%$ of comments vs. betrayal: $1.2\%$). This suggests users maintain generally positive relationships with AI systems, but trust appears fragile and concentrated around specific trigger events. Betrayal language clustered significantly around model update periods ($\chi^2=23.47$, $p<0.001$), indicating that trust erosion is often precipitated by perceived capability losses rather than absolute performance metrics.

\subsection{Eleven Distinct Mutual Wanting User Types}

K-means clustering with silhouette optimization identified eleven distinct user types based on mutual wanting patterns (optimal $K=11$, silhouette score $=0.304$). Table~\ref{tab:user_clusters} presents the distribution and characteristics of these clusters.

\begin{table}[htbp]
\centering
\caption{Mutual Wanting User Type Distribution and Characteristics}
\label{tab:user_clusters}
\begin{tabular}{@{}llrp{5cm}@{}}
\toprule
Cluster & User Type & \% & Key Characteristics \\
\midrule
C0 & Stable Users & 18.3\% & High reliability wanting, minimal tension \\
C1 & Attached Users & 14.7\% & High anthropomorphism, expectation violations \\
C2 & Power Users & 12.9\% & System efficiency focus, professional language \\
C3 & Creative Users & 11.4\% & Creativity wanting, loss language patterns \\
C4 & Evaluator Users & 10.8\% & Intelligence focus, model comparisons \\
C5 & Collaborative Users & 9.6\% & Helpfulness wanting, active feedback \\
C6 & Resistant Users & 8.7\% & High tension, change resistance \\
C7 & Satisfied Users & 6.9\% & Balanced patterns, trust language \\
C8 & Emotional Users & 6.1\% & Emotional language, relationship terms \\
C9 & Technical Users & 5.4\% & Technical focus, efficiency optimization \\
C10 & Confused Users & 4.2\% & Mixed patterns, uncertainty language \\
\bottomrule
\end{tabular}
\end{table}

Each cluster showed distinct communication patterns and response preferences, suggesting the need for personalized interaction strategies rather than one-size-fits-all approaches.

\subsection{Expectation Violation Patterns}

Our expectation analysis identified measurable patterns of user disappointment and violated expectations. Explicit expectation violations appeared in 2.23\% of comments (499 instances), clustering around specific linguistic patterns:

\begin{itemize}
    \item "Not what I expected" ($234$ instances)
    \item "Used to work better" ($187$ instances)  
    \item "Thought it would be different" ($156$ instances)
\end{itemize}

These violations were not randomly distributed but showed significant correlation with model update periods and specific capability domains (performance, creativity, personality traits).

\subsection{GPT-5 Release Impact Analysis}

The GPT-5 release provided a natural experiment for measuring mutual wanting dynamics during major model transitions. Table~\ref{tab:gpt5_impact} summarizes the key changes observed.

\begin{table}[htbp]
\centering
\caption{GPT-5 Release Impact on User Sentiment and Concerns}
\label{tab:gpt5_impact}
\begin{tabular}{@{}lrrr@{}}
\toprule
Metric & Pre-GPT-5 & Post-GPT-5 & Change \\
\midrule
\textbf{Sentiment Metrics} & & & \\
Compound Score & $0.479$ & $0.435$ & $-0.044$* \\
Anger Rate & $0.001$ & $0.002$ & $+38.18\%$ \\
Joy Rate & $0.002$ & $0.002$ & $-6.65\%$ \\
\midrule
\textbf{User Concerns (\%)} & & & \\
Performance & $11.0\%$ & $13.0\%$ & $+2.02$pp \\
Safety & $6.6\%$ & $8.6\%$ & $+1.94$pp \\
Accuracy & $18.0\%$ & $18.9\%$ & $+0.87$pp \\
Capabilities & $20.1\%$ & $18.9\%$ & $-1.20$pp \\
\midrule
\textbf{Expectation Dynamics} & & & \\
Expectation Comments & $134$ & - & - \\
Reality Comments & - & $3,369$ & - \\
Expectation-Reality Gap & - & - & $-0.27$ \\
\bottomrule
\end{tabular}
\begin{tablenotes}
\small
\item[*] Statistically significant at $p<0.05$
\end{tablenotes}
\end{table}

\subsubsection{Sentiment and Emotional Shifts}

Overall sentiment became significantly more negative following GPT-5's release (compound score change: $-0.0441$, $p=0.0312$). Anger increased by $38.18\%$, while joy decreased by $6.65\%$. The expectation-reality gap measured $-0.27$, indicating that user reality fell substantially short of pre-release expectations.

\subsubsection{Concern Pattern Changes}

User concerns shifted significantly post-release:
\begin{itemize}
    \item Performance concerns: $+2.02$ percentage points ($+18.4\%$ relative increase)
    \item Safety concerns: $+1.94$ percentage points ($+29.4\%$ relative increase)  
    \item Accuracy concerns: $+0.87$ percentage points ($+4.9\%$ relative increase)
    \item Capability concerns: $-1.20$ percentage points ($-6.0\%$ relative decrease)
\end{itemize}

\subsection{API Probe Model Persona Analysis}

Controlled API probing revealed distinct persona characteristics across the 9 tested models. Response patterns varied significantly across dimensions of warmth, formality, and response length. Table~\ref{tab:model_personas} summarizes key persona metrics across models.

\begin{table}[htbp]
\centering
\caption{Model Persona Characteristics from API Probe Analysis}
\label{tab:model_personas}
\begin{tabular}{@{}lrrr@{}}
\toprule
Model & Avg Length & Warmth Score & Formality Score \\
\midrule
gpt-3.5-turbo & 804 & 0.14 & 0.11 \\
gpt-4 & 898 & 0.09 & 0.22 \\
gpt-4o & 1018 & 0.07 & 0.05 \\
gpt-4.1 & 907 & 0.05 & 0.04 \\
o3 & 363 & 0.11 & 0.02 \\
gpt-4.1-mini & 846 & 0.19 & -0.06 \\
gpt-4o-mini & 947 & 0.09 & 0.01 \\
gpt-5 & 8 & 0.00 & 0.00 \\
gpt-5-mini & 45 & 0.00 & 0.00 \\
\bottomrule
\end{tabular}
\end{table}

Notable patterns include gpt-3.5-turbo showing the highest warmth scores, gpt-4 exhibiting the highest formality, and both GPT-5 variants showing dramatically reduced response lengths with zero warmth/formality scores. These differences align with user perceptions of personality changes, providing objective validation of subjective user reports.

\subsection{Topic Modeling and Discourse Themes}

Dual-algorithm topic modeling (LDA + NMF) revealed convergent themes across both approaches. The most prominent topics were Performance Complaints (weight=0.089), Personality Changes (weight=0.078), Feature Requests (weight=0.071), Model Comparisons (weight=0.067), and Trust \& Reliability (weight=0.063). Performance complaints showed the largest increase post-GPT-5 release ($\Delta$weight=+0.024), consistent with our concern analysis.

\section{Discussion}

\subsection{Implications for AI System Design}

Our findings have profound implications for how AI systems should be designed and deployed. The $48.65\%$ anthropomorphism rate suggests that human-like attribution is not a design bug but a fundamental human response that requires accommodation. Rather than discouraging anthropomorphization, systems should be designed to safely support these attributions while maintaining appropriate boundaries.

The identification of 11 distinct user types challenges one-size-fits-all approaches to AI interaction. "Stable Users" prioritizing reliability may require different communication patterns than "Creative Users" mourning lost capabilities or "Technical Users" seeking efficiency optimization. This suggests the need for adaptive systems that can recognize and respond to different mutual wanting profiles.

\subsection{Trust as a Fragile Asset}

The $11.6:1$ trust-to-betrayal ratio indicates that users maintain generally positive relationships with AI systems, but this trust appears fragile. The concentration of betrayal language around model update periods suggests that trust erosion is often triggered by perceived personality changes rather than absolute performance metrics. This highlights the importance of managing not just technical capabilities but relational continuity during system updates.

\subsection{Expectation Violation as Early Warning System}

The measurable patterns of expectation violations (2.23\% of discourse) provide a potential early warning system for user dissatisfaction. The clustering of violations around specific linguistic patterns ("not what I expected," "used to work better") enables automated monitoring systems that could detect and address user concerns before they escalate to community-wide discussions.

\subsection{The Paradox of Mutual Wanting}

Our results reveal a fundamental tension in mutual wanting dynamics: users want AI systems to be reliable, consistent, and trustworthy, while simultaneously expecting continuous improvement and capability expansion. AI systems, through their optimization objectives, "want" clear inputs and structured interactions, but must balance this with user desires for natural, relationship-like communication. This paradox suggests that successful AI development requires explicit management of competing wants rather than optimizing for single objectives.

\section{Limitations and Future Work}

Our analysis has several limitations that future work should address. The Reddit-based dataset, while large and naturalistic, may not represent all user populations. More casual users who don't engage in online discourse might have different mutual wanting patterns. Additionally, our temporal analysis focuses on a single major model transition (GPT-5 release); patterns might vary for different types of updates or AI systems.

Future work should expand this analysis across multiple platforms, cultural contexts, and model architectures. Longitudinal studies tracking individual users across multiple model transitions could provide insights into adaptation patterns and long-term relationship dynamics. Intervention studies testing different communication strategies during model updates could validate our recommendations for expectation management.

The methodology could be enhanced through cross-platform validation, inclusion of non-English discourse, and integration with objective performance metrics. Future research might also explore causal relationships between specific design changes and user response patterns.

\section{Conclusion}

This work provides the first large-scale empirical validation of mutual wanting dynamics in human-AI interaction. Our analysis of 22,411 user comments and 729 controlled API responses reveals that mutual wanting is not just a theoretical concept but a measurable phenomenon with clear patterns and implications.

The identification of $48.65\%$ anthropomorphism rates, $11.6:1$ trust-betrayal ratios, and $11$ distinct user types provides concrete targets for AI system design. The development of expectation violation detection capabilities and trust monitoring systems offers practical tools for managing human-AI relationships during the rapid pace of AI development.

Most importantly, our findings suggest that the future of AI development cannot ignore the relational dimension of human-AI interaction. As AI systems become more sophisticated and ubiquitous, understanding and aligning mutual wants becomes not just a research curiosity but a practical necessity for building trustworthy and sustainable AI systems.

The Mutual Wanting Alignment Framework (M-WAF) provides a foundation for this work, but much remains to be done. By taking seriously the bidirectional nature of human-AI desires, we can move toward AI systems that are not just more capable, but more relationally aware and aligned with human values.

\bibliographystyle{plain}
\bibliography{references}

\newpage

\section*{Agents4Science AI Involvement Checklist}

\begin{enumerate}
    \item \textbf{Hypothesis development}: Hypothesis development includes the process by which you came to explore this research topic and research question. This can involve the background research performed by either researchers or by AI. This can also involve whether the idea was proposed by researchers or by AI. 

    Answer: \involvementC{} 
    
    Explanation: The core "mutual wanting" hypothesis emerged from AI analysis of user discourse patterns, but was developed through extensive human-AI collaboration. AI identified initial patterns in anthropomorphic language and trust dynamics, while humans provided theoretical framing and research question refinement. The bidirectional concept was jointly developed through iterative discussion and validation.

    \item \textbf{Experimental design and implementation}: This category includes design of experiments that are used to test the hypotheses, coding and implementation of computational methods, and the execution of these experiments. 

    Answer: \involvementC{} 
    
    Explanation: AI conducted the majority of coding implementation, including the 47-dimensional feature extraction pipeline, clustering algorithms, and statistical analysis. However, humans designed the overall experimental framework, selected methodological approaches, and provided oversight for data collection protocols. The API probe design and validation were collaborative efforts.

    \item \textbf{Analysis of data and interpretation of results}: This category encompasses any process to organize and process data for the experiments in the paper. It also includes interpretations of the results of the study.
 
    Answer: \involvementC{} 
    
    Explanation: AI performed most of the data processing, statistical analysis, and pattern identification from the 22,411 comments and 729 API responses. However, interpretation of results, particularly the sociological implications and practical applications, was primarily human-driven. The discovery of 11 user types was AI-identified but human-interpreted for meaning and implications.

    \item \textbf{Writing}: This includes any processes for compiling results, methods, etc. into the final paper form. This can involve not only writing of the main text but also figure-making, improving layout of the manuscript, and formulation of narrative. 

    Answer: \involvementC{} 
    
    Explanation: AI generated the majority of the paper text, structured arguments, and created tables/figures. However, humans provided significant guidance on narrative framing, edited for clarity and academic style, and ensured proper alignment with conference requirements. The overall storytelling and presentation strategy was collaboratively developed.

    \item \textbf{Observed AI Limitations}: What limitations have you found when using AI as a partner or lead author? 

    Description: AI occasionally over-interpreted statistical patterns without sufficient theoretical grounding and sometimes lacked nuanced understanding of sociological context. AI required human oversight for ethical considerations and potential biases in anthropomorphism analysis. Additionally, AI needed guidance on research ethics and appropriate claims strength given data limitations.
\end{enumerate}

\newpage

\section*{Agents4Science Paper Checklist}

\begin{enumerate}

\item {\bf Claims}
    \item[] Question: Do the main claims made in the abstract and introduction accurately reflect the paper's contributions and scope?
    \item[] Answer: \answerYes{} 
    \item[] Justification: The abstract and introduction clearly state our empirical findings (48.65\% anthropomorphism, 11 user types) and scope (Reddit discourse + API probing). Claims are supported by results sections.

\item {\bf Limitations}
    \item[] Question: Does the paper discuss the limitations of the work performed by the authors?
    \item[] Answer: \answerYes{} 
    \item[] Justification: Section 6 explicitly discusses limitations including Reddit-only population, single model transition focus, and potential platform bias.

\item {\bf Theory assumptions and proofs}
    \item[] Question: For each theoretical result, does the paper provide the full set of assumptions and a complete (and correct) proof?
    \item[] Answer: \answerNA{} 
    \item[] Justification: This is an empirical study without formal theoretical proofs. Our framework is empirically validated rather than theoretically proven.

\item {\bf Experimental result reproducibility}
    \item[] Question: Does the paper fully disclose all the information needed to reproduce the main experimental results of the paper to the extent that it affects the main claims and/or conclusions of the paper (regardless of whether the code and data are provided or not)?
    \item[] Answer: \answerYes{} 
    \item[] Justification: Methodology section provides detailed parameters ($K=11$ clustering, $47$-dimensional features, specific statistical tests) and data collection procedures.

\item {\bf Open access to data and code}
    \item[] Question: Does the paper provide open access to the data and code, with sufficient instructions to faithfully reproduce the main experimental results, as described in supplemental material?
    \item[] Answer: \answerNo{} 
    \item[] Justification: Reddit data contains potentially sensitive user comments requiring privacy protection. API probe data and analysis code could be made available with appropriate anonymization.

\item {\bf Experimental setting/details}
    \item[] Question: Does the paper specify all the training and test details (e.g., data splits, hyperparameters, how they were chosen, type of optimizer, etc.) necessary to understand the results?
    \item[] Answer: \answerYes{} 
    \item[] Justification: Methodology section specifies clustering parameters ($K=3$-$15$ optimization), topic modeling setup ($15$ topics each for LDA/NMF), and statistical testing procedures.

\item {\bf Experiment statistical significance}
    \item[] Question: Does the paper report error bars suitably and correctly defined or other appropriate information about the statistical significance of the experiments?
    \item[] Answer: \answerYes{} 
    \item[] Justification: Results include p-values ($p=0.0312$ for sentiment changes), chi-square statistics ($\chi^2=23.47$), and effect sizes throughout the analysis.

\item {\bf Experiments compute resources}
    \item[] Question: For each experiment, does the paper provide sufficient information on the computer resources (type of compute workers, memory, time of execution) needed to reproduce the experiments?
    \item[] Answer: \answerNo{} 
    \item[] Justification: We focused on methodological details over computational requirements. Future versions should include resource specifications for NLP processing and clustering analysis.

\item {\bf Code of ethics}
    \item[] Question: Does the research conducted in the paper conform, in every respect, with the Agents4Science Code of Ethics (see conference website)?
    \item[] Answer: \answerYes{} 
    \item[] Justification: Research uses publicly available data with appropriate privacy considerations and focuses on beneficial applications for AI system improvement.

\item {\bf Broader impacts}
    \item[] Question: Does the paper discuss both potential positive societal impacts and negative societal impacts of the work performed?
    \item[] Answer: \answerYes{} 
    \item[] Justification: Discussion section addresses positive impacts (better AI systems, trust management) and implicitly addresses risks through emphasis on ethical anthropomorphism design.

\end{enumerate}

\end{document}
